% ============================================
% Johns Hopkins Letter Template
% Assistant Professor Cover Letter – La Salle University
% ============================================

\documentclass[11pt,letterpaper]{letter}

\usepackage{graphicx}
\usepackage{eso-pic}
\usepackage{geometry}
\usepackage{microtype}
\usepackage[hidelinks]{hyperref}

% ----- Page Setup -----
\geometry{left=25mm,right=25mm,top=25mm,bottom=25mm}

% ----- Logo Placement (first page only) -----
\newcommand{\PutJHULogoFG}[1]{%
  \AddToShipoutPictureFG*{%
    \AtPageUpperLeft{%
      \hspace{10mm}%
      \raisebox{-38mm}{%
        \includegraphics[width=6cm]{#1}%
      }%
    }%
  }%
}

% ----- Sender -----
\signature{Decory Edwards}
\address{Department of Economics \\ Johns Hopkins University \\ Baltimore, MD 21218 \\ \texttt{dedwar65@jh.edu}}

\begin{document}
\PutJHULogoFG{Images/jhulogo.png}

\begin{letter}{Search Committee \\
Department of Humanity and Society \\
La Salle University \\
Philadelphia, PA 19141 \\ USA}

\opening{Dear Members of the Search Committee,}

I am writing to express my interest in the tenure-track Assistant Professor position in Economics in the Department of Humanity and Society at La Salle University, beginning Fall 2026. I am a Ph.D. candidate in Economics at Johns Hopkins University, advised by Professors Christopher D.~Carroll, Nicholas W.~Papageorge, and Stelios Fourakis, and I expect to receive my degree in May 2026. My research fields are macroeconomics, wealth and income dynamics, and computational economics.

My research agenda focuses on understanding the determinants of wealth inequality and the mechanisms through which heterogeneity in returns to wealth shapes aggregate outcomes. In my job-market paper, I estimate the degree of heterogeneity in individual portfolio returns using micro-data from the Survey of Consumer Finances and simulate a structural model in which households face idiosyncratic return risk. I show that allowing for persistent heterogeneity in returns substantially improves the model’s ability to match the empirical distribution of wealth, offering a tractable way to connect micro-level return dispersion to macroeconomic inequality.

In a second project, I use the Health and Retirement Study to examine the relationship between interpersonal trust and portfolio performance. Building on the literature that finds a hump-shaped relationship between trust and income, I show that a similar relationship holds for individual returns to wealth measured in the HRS data, highlighting a behavioral channel through which social attitudes shape saving and investment outcomes. This work also provides new estimates of the persistent component of returns to net worth, which motivate the calibration of heterogeneity in my structural model.

I am strongly drawn to La Salle University’s mission as a Catholic Lasallian institution devoted to teaching excellence, student success, and the formation of engaged citizens. The department’s commitment to undergraduate education, inclusive pedagogy, and a supportive intellectual community aligns well with my teaching philosophy and desire to mentor students both inside and outside the classroom. I welcome the opportunity to contribute to the School of Arts \& Sciences by teaching \textit{Principles of Microeconomics} and \textit{Principles of Macroeconomics}, in addition to upper-level courses in economics, statistics, and econometrics. I also look forward to advising majors and participating in service to the department and university, including evening courses and weekend events when needed.

\textbf{Teaching.} My teaching experience centers on undergraduate macroeconomics and an upper-level macro policy seminar, where I have led weekly discussion sections and held regular office hours for classes ranging from 16 to 40 students. My approach is student-centered: I aim to help students “convince themselves” that they have moved from confusion to understanding by holding structured office-hour coaching that builds trust and confidence. At La Salle, I am prepared to teach a range of courses including \textit{Principles of Microeconomics}, \textit{Principles of Macroeconomics}, \textit{Statistics}, \textit{Econometrics}, and upper-level undergraduate electives aligned with my research. I also welcome the opportunity to mentor students in undergraduate research and support them as they develop analytical, quantitative, and critical reasoning skills.

I believe I would be a strong fit because my empirical and computational training allows me to (i) deliver rigorous instruction in core and advanced economics courses, (ii) integrate data-driven methods and reproducible workflows into the curriculum, and (iii) contribute to La Salle’s student-focused mission by fostering inclusive and supportive learning environments. I am enthusiastic about joining a university community that values teaching excellence, collaboration, and the holistic development of students.

I would be honored to join the Department of Humanity and Society at La Salle University. Thank you for your consideration; I welcome the opportunity to discuss my fit with your department at your convenience.

\closing{Sincerely,}

\end{letter}
\end{document}
